\scanpage{72}
\uline{Remarque} Dans deux contextes différents
``deploiement''$^{*}$ de \deleted{\unsure{strag}} \unsure{stratification|stratifications}, on
a rencontré des espaces topologiques ordonnés\inserted{$\Sigma$},
satisfaisant les conditions suivantes

\margin{$*$ en deux sens assez différents}

a) $\Sigma_{(1)}\xrightarrow{\uline{s}}\Sigma\deleted{\unsure{0|*}}=\Sigma(0)$
est une fibration propre
% Detached ink point just above-left of the equality, approximately (776,495);
% another detached point between the a and b rows, approximately (812,562).
% Their function is unassigned.

b) $\Sigma(1)\xrightarrow{\uline{b}}\Sigma(0)$ est une immersion propre

c) Si $x,y\in\Sigma$ \deleted{\unsure{sont}} \unsure{ds|dans} une $\unsure{\overline{\mathrm m}|=}$ composante connexe,
alors $x\leq y\Longrightarrow x=y$
\inserted{(i.e. si $x,y\in\Sigma$, $x<y$, alors $x,y$ sont dans deux
\unsure{composantes connexes différentes})}
% The terminal inserted words overlap the following baseline and meet the edge.
% No independent extra cancelled group has been isolated there.

\margin{conséquence de a) b) si on suppose que
$\Sigma(1)^{*}\xrightarrow{\uline b}\Sigma$ est \uline{\unsure{rare}},
et que les comp. connexes de $\Sigma$ de dim finie...}

Dans un cas (celui d'une notion assez \unsure{anodine}\unsure{)}
\unsure{de la|la} notion de déploiement d'une stratification de type $I$ disons,
par $\Sigma=\coprod X_i$,
$\Sigma_1=\displaystyle\coprod_{\substack{i,j\text{ tels que}\\X_i\subset X_j}}X_i$),
$\uline s$ est étale, (\unsure{localt} fini donc) et rien de plus spécial pour $\uline b$.
Dans le cas de la théorie \deleted{\unsure{présente|proposée}}\inserted{développée}
ici des déploiements\inserted{de stratifications équisingulières},
$\uline b$ est bordant, rien de spécial sur $\uline s$.
Dans l'un et l'autre cas, on peut appliquer le sorite
du $\widetilde I=\pi_0(\Sigma)$ --- dans le premier cas,
% Two small detached points occur below-left of the closing parenthesis,
% approximately (838,1842) and (834,1851); their function is unassigned.
cela revient à remplacer la stratification par les $X_i$ par celle
par les comp. connexes des $X_i$ (donnant un $\widetilde I$;
dans le deuxième,
\deleted{à \unsure{remp|rempl} \unsure{ito|itou} $\unsure{\Longleftarrow|=}$ \unsure{exact}}
itou pour les
\[
X_i^{*}=X_i\setminus\bigcup_{j<i}X_j.
\]
\margin{\handwritten{71}}
